Sources : études citées, présentées dans le complément « Les élasticités estimées sur données françaises » ; chaque chiffre, avec sa page et son tableau, figure dans donnees/elasticites\_france.csv.\par Lecture : chaque panneau compare des estimations d'une même étude. Les quatre mesurent la même grandeur et partagent la même échelle, mais portent sur des revenus, des années et des populations différents. Cabannes et al. ne publient pas d'écart-type : seul le décile supérieur est significatif, et un autre de leurs modèles y donne 0,50. Chez Lehmann et al., les réformes qui identifient l'effet ne touchaient que les salariés gagnant moins de deux fois le salaire minimum.
